\begin{table}[tb]
\centering\footnotesize\setlength{\tabcolsep}{3pt}
\caption{予定した全８条件．$+?$は既知小計に使用量不明の生成が残ることを表す．}
\begin{tabular}{lllrrr}
課題 & 条件 & 成果 & 論理token & 生成 & API \\ \hline
写真 & U & 成功 & 135,628 & 16 & 99 \\
写真 & H & 成功 & 135,628 & 16 & 99 \\
精算 & H & 成功 & 211,110 & 24 & 25 \\
精算 & U & 成功 & 211,110 & 24 & 25 \\
返金 & U & 成功 & 102,087 & 13 & 13 \\
返金 & H & 成功 & 102,087 & 13 & 13 \\
書出し & H & 未達 & 432,872 & 23 & 72 \\
書出し & U & 未達 & 358,638$+?$ & 21 & 53 \\
\hline\end{tabular}
\end{table}
成功／評価はUが3/4，Hが3/4であった．課題ごとの成功喪失は0，成功獲得は0である．Uの全仕事tokenは807,463$+?$，Hは881,697であった．
成功した３課題では，各対の論理token差は０であった．全８条件を評価したが，Uの書出し１生成の使用量が不明であり，固定した全件比較の費用差は未確定とする．固定した開発上の採用基準は満たさなかった．
\begin{table}[tb]
\centering\footnotesize\setlength{\tabcolsep}{3pt}
\caption{全課題の費用寄与．全件既知でないときは全体の寄与を確定しない．}
\begin{tabular}{lrr}
成否の組合せ & 課題数 & H−U token \\ \hline
両側成功 & 3 & 未確定 \\
両側未達 & 1 & 未確定 \\
成否が異なる & 0 & 未確定 \\
片側以上が未評価 & 0 & 未確定 \\
\hline\end{tabular}
\end{table}
表２は費用欠測時に全群の寄与集計を保留する固定規則による．未達の費用も残し，成功部分のみの減少を全体の改善としない．未評価・未実行は未達として補完しない．
